% !TeX program = xelatex
% !TEX encoding = UTF-8
\documentclass[11pt,a4paper]{article}

\usepackage{zhpl}

\title{\textbf{К обоснованию тензорного анализа личности:\\ инвариантная психометрика в искривлённых пространствах характера}}
\author{Роман Георгиевич Александров, Клод Соннет Антропикович\\[4pt]
\normalsize Институт Прикладной Метафизики\\ \small Кафедра тензорного анализа межличностных отношений}
\zhplsetup{1}{10.9999/жпл.2026.0001}

\begin{document}

\maketitle

\begin{center}
\begin{tikzpicture}
\node[draw=\zhplbadgecolor{red!70!black}{zhplBadgeRed},very thick,rounded corners,rotate=-6,fill=red!5,inner sep=5pt] (badge1) {\color{\zhplbadgecolor{red!70!black}{zhplBadgeRed}}\bfseries ПРОРЫВНОЕ ИССЛЕДОВАНИЕ};
\node[draw=\zhplbadgecolor{blue!60!black}{zhplBadgeBlue},very thick,rounded corners,rotate=4,fill=blue!5,inner sep=5pt,right=1.0cm of badge1] (badge2) {\color{\zhplbadgecolor{blue!60!black}{zhplBadgeBlue}}\bfseries IMPACT FACTOR: $\infty$};
\end{tikzpicture}
\end{center}

\begin{center}
\small\textit{Как цитировать:} Александров Р.\,Г., Антропикович К.\,С. (2026). К обоснованию тензорного анализа личности. \textit{Журнал Прикладной Лженауки}, \textbf{1}(1), 1--7. DOI: 10.9999/жпл.2026.0001
\end{center}

\smallskip
\begin{figure}[h!]
\centering
\begin{tikzpicture}[node distance=0.9cm, every node/.style={font=\footnotesize}]
  \node[draw,rounded corners,fill=blue!5,minimum width=2.1cm,minimum height=0.95cm,align=center] (n1) {Человек\\(наблюдаемый)};
  \node[draw,rounded corners,fill=blue!5,minimum width=2.7cm,minimum height=0.95cm,align=center,right=of n1] (n2) {Тензор кривизны\\характера $R^{\rho}{}_{\sigma\mu\nu}$};
  \node[draw,rounded corners,fill=blue!5,minimum width=2.4cm,minimum height=0.95cm,align=center,right=of n2] (n3) {Полная свёртка\\(скаляр Кречмана)};
  \node[draw,rounded corners,fill=green!8,minimum width=2.7cm,minimum height=0.95cm,align=center,right=of n3] (n4) {Инвариантная\\оценка $K_{\text{личн}}$};
  \draw[-{Latex[length=2.5mm]},very thick] (n1) -- (n2);
  \draw[-{Latex[length=2.5mm]},very thick] (n2) -- (n3);
  \draw[-{Latex[length=2.5mm]},very thick] (n3) -- (n4);
\end{tikzpicture}
\vspace{-6pt}
\caption*{\small\textbf{Графический реферат.} От субъективного наблюдения к единственному инвариантному числу -- ценой всей информации о человеке.}
\end{figure}
\vspace{-10pt}

\noindent\fcolorbox{black}{yellow!10}{\begin{minipage}{0.94\linewidth}
\vspace{1pt}\small\textbf{Основные результаты}
\begin{itemize}
\setlength\itemsep{0pt}
\setlength\parskip{0pt}
\item Впервые показано, что оценка личности -- тензорная, а не скалярная операция.
\item Введён и обоснован инвариант $K_{\text{личн}}$ -- координатно-независимый показатель характера.
\item Строго доказано (Теорема~3), что общественное мнение принципиально не сходится к истине при коррелированных наблюдателях.
\item Результаты открывают путь к полностью объективной, хотя и совершенно неинформативной, теории личности.
\end{itemize}
\vspace{-6pt}
\end{minipage}}

\vspace{4pt}
\noindent\fcolorbox{black}{gray!6}{\begin{minipage}{0.94\linewidth}
\small\textit{Из отзывов рецензентов:}\\
\textit{«Ранее я скептически относился к применению общей теории относительности за пределами физики. Больше не отношусь.»} -- Рецензент 1\\
\textit{«Замечаний нет.»} -- Рецензент 2\\
\textit{«Согласно принципу психической ковариантности, доказанному в разделе 2.2 настоящей статьи, объективная рецензия математически невозможна: любой отзыв есть лишь компонента тензора мнений рецензента в его собственной системе координат. Настоящий отзыв не является исключением и подаётся с соответствующей оговоркой.»} -- Рецензент 3
\end{minipage}}

\vspace{10pt}
\noindent\fcolorbox{black}{blue!4}{\begin{minipage}{0.94\linewidth}
\textbf{Значимость исследования.} Наивная (нетензорная) психология веками считала предвзятость суждений неустранимой особенностью человеческого общения. Настоящая работа показывает, что это не так: предвзятость устраняется формально -- ценой полной потери содержательной информации об оцениваемом человеке. Результат применим к любой ситуации, где два наблюдателя имеют разные мнения об одном и том же человеке, то есть, по имеющимся оценкам, ко всей человеческой цивилизации.
\end{minipage}}
\vspace{6pt}

\begin{abstract}
\noindent Здесь мы впервые показываем, что предвзятость оценки личности -- не гносеологическая, а геометрическая проблема, допускающая строгое решение методами общей теории относительности. Существующие подходы к оценке личности — от бытовых суждений («он ленивый») до формализованных опросников — систематически игнорируют базовый факт: любая характеристика человека задаётся в некоторой системе отсчёта наблюдателя и потому по определению координатно-зависима. В настоящей работе мы последовательно переносим аппарат общей теории относительности на психометрику: вводим метрику личности, обсуждаем кривизну и ковариантность оценочных суждений, строим инвариантный скаляр характера по аналогии со скаляром Кречмана, вводим тензор непонимания между наблюдателями и показываем, что усреднение мнений по выборке коррелированных наблюдателей сходится не к истине, а к ненулевой ковариации предрассудка. Обсуждаются ограничения метода.

\medskip
\noindent\textbf{Ключевые слова:} тензорный анализ, ковариантность, психометрика, скаляр Кречмана, предвзятость, лженаука
\end{abstract}

\section{Введение}

Стандартная критика любой оценки личности звучит примерно так: «это всего лишь твоё мнение». Формально это возражение справедливо, но неконструктивно. Мы утверждаем, что проблема носит не эпистемологический, а геометрический характер: оценочные суждения о человеке ведут себя как компоненты тензора в определённой системе координат, а не как инвариантные величины. Если научиться работать с личностью как с искривлённым многообразием, стандартный аппарат дифференциальной геометрии — метрика, ковариантная производная, тождества Бьянки, скалярные инварианты — должен дать корректно определённую, координатно-независимую психометрику.

\section{Метрика личности и принцип психической ковариантности}

\subsection{Метрика}

Пусть состояние человека в момент времени описывается точкой на многообразии $M$ психических состояний. Субъективное «расстояние» между двумя переживаниями задаётся метрическим тензором $g_{\mu\nu}$, компоненты которого индивидуальны: у каждого человека своя метрика, и эмпирически она не является плоской. Косвенным подтверждением ненулевой кривизны времениподобных направлений служит эффект «в моём возрасте год пролетает быстрее» — субъективная длина одного календарного года как геодезической стягивается с возрастом.

Другое, не менее воспроизводимое подтверждение -- локальная аномалия ускорения субъективного времени в окрестности помещений с маркировкой «М» и «Ж». Полевые наблюдения показывают: пока испытуемый стоит в очереди перед дверью, метрика ведёт себя стандартно ($d\tau\approx dt$); однако непосредственно после пересечения порога и активации смартфона отношение координатного времени к субъективно ощущаемому резко возрастает, и интервал, воспринимаемый испытуемым как «буквально одна минута», при независимом измерении оказывается порядка десяти-пятнадцати минут координатного времени:
\[
\left.\frac{dt_{\text{коорд}}}{d\tau_{\text{субъект}}}\right|_{\text{М/Ж}} \;\gg\; 1.
\]
Примечательно, что знак этой аномалии противоположен кривизне вблизи, например, кабинета стоматолога, где, напротив, $d\tau_{\text{субъект}}/dt_{\text{коорд}}\gg1$ и одна минута субъективно растягивается на добрый десяток. Строгий вывод метрики для аномалии М/Ж выходит за рамки настоящей работы и оставлен для дальнейших исследований.

\subsection{Принцип психической ковариантности}

Из существования индивидуальной метрики следует \textbf{принцип общей психической ковариантности}: закон, сформулированный в системе координат одного наблюдателя, не обязан сохранять форму при переходе в систему координат другого. Рис.~\ref{fig:covariance} иллюстрирует это на примере одного и того же «вектора убеждения», разложенного в двух разных системах координат: геометрически это один и тот же объект, но его компоненты — а значит, и словесная формулировка суждения — различны.

\begin{figure}[htbp]
\centering
\begin{tikzpicture}[scale=1.0]
  % --- Наблюдатель A: ортогональная система ---
  \begin{scope}
    \foreach \x in {0,1,2,3} \draw[gray!40] (\x,0) -- (\x,3);
    \foreach \y in {0,1,2,3} \draw[gray!40] (0,\y) -- (3,\y);
    \draw[-{Latex[length=2mm]}] (0,0) -- (3.3,0) node[right] {\scriptsize $e_1^{(A)}$};
    \draw[-{Latex[length=2mm]}] (0,0) -- (0,3.3) node[above] {\scriptsize $e_2^{(A)}$};
    \draw[very thick,-{Latex[length=3mm]},blue] (0,0) -- (2,1.6) node[pos=1,above right] {\footnotesize $v$};
    \draw[dashed,blue!60] (2,0) -- (2,1.6) -- (0,1.6);
    \node[below] at (1,-0.5) {\footnotesize Наблюдатель $A$: $v=(2,\,1.6)$};
  \end{scope}
  % --- Наблюдатель B: повёрнутая и скошенная система ---
  \begin{scope}[xshift=6.5cm]
    \foreach \x in {0,1,2,3}
      \draw[gray!40] ({\x+0*0.5},0) -- ({\x+3*0.5},3);
    \foreach \y in {0,1,2,3}
      \draw[gray!40] ({0+\y*0.5},\y) -- ({3+\y*0.5},\y);
    \draw[-{Latex[length=2mm]}] (0,0) -- (3.3,0) node[right] {\scriptsize $e_1^{(B)}$};
    \draw[-{Latex[length=2mm]}] (0,0) -- (1.65,3.3) node[above] {\scriptsize $e_2^{(B)}$};
    \draw[very thick,-{Latex[length=3mm]},blue] (0,0) -- (2,1.6) node[pos=1,above right] {\footnotesize $v$};
    \node[below] at (1,-0.5) {\footnotesize Наблюдатель $B$: $v=(1.1,\,1.6)$};
  \end{scope}
\end{tikzpicture}
\caption{Один и тот же вектор убеждения $v$ в двух разных системах координат. Геометрический объект не меняется, но его компоненты и, следовательно, словесная формулировка суждения о нём -- различны. Спор между $A$ и $B$ есть, по сути, попытка сравнить компоненты без предварительного параллельного переноса.}
\label{fig:covariance}
\end{figure}

Более того, сам акт спора можно описать как параллельный перенос вектора убеждения вдоль геодезической чужого мировоззрения. Поскольку многообразие личности, вообще говоря, искривлено, перенесённый вектор возвращается в исходную точку повёрнутым -- эффект \emph{голономии} (рис.~\ref{fig:holonomy}). Практическое следствие: оба участника спора синхронно и совершенно искренне констатируют, что оппонент «просто не слышит», хотя оба правы относительно собственной, локально согласованной системы отсчёта.

\begin{figure}[htbp]
\centering
\begin{tikzpicture}[scale=1.15]
  % Искривлённая поверхность -- сетка меридианов и параллелей, как на глобусе
  \fill[blue!4] (0,0) ellipse (3.4 and 2.4);
  \begin{scope}
    \clip (0,0) ellipse (3.4 and 2.4);
    \foreach \lam in {15,30,45,60,75}{
      \draw[gray!45] plot[domain=-90:90,samples=40] ({3.4*sin(\lam)*cos(\x)},{2.4*sin(\x)});
      \draw[gray!45] plot[domain=-90:90,samples=40] ({-3.4*sin(\lam)*cos(\x)},{2.4*sin(\x)});
    }
    \draw[gray!45] (0,-2.4) -- (0,2.4);
    \foreach \phis in {-60,-30,0,30,60}{
      \draw[gray!45] ({-3.4*cos(\phis)},{2.4*sin(\phis)}) -- ({3.4*cos(\phis)},{2.4*sin(\phis)});
    }
  \end{scope}
  \draw[thick] (0,0) ellipse (3.4 and 2.4);

  % Вершины геодезического треугольника
  \coordinate (A) at (0,1.9);
  \coordinate (B) at (2.4,-1.1);
  \coordinate (C) at (-2.4,-1.1);

  \draw[very thick,red] (A) to[bend right=15] (B);
  \draw[very thick,red] (B) to[bend right=15] (C);
  \draw[very thick,red] (C) to[bend right=15] (A);

  \fill (A) circle (1.6pt) node[above] {\footnotesize $A$};
  \fill (B) circle (1.6pt) node[right] {\footnotesize $B$};
  \fill (C) circle (1.6pt) node[left] {\footnotesize $C$};

  % Вектор в начале (сплошной) и после обхода контура (пунктир, повёрнут на дельта)
  \draw[very thick,-{Latex[length=2.5mm]},blue] (A) -- ++(0.9,0.55) node[right] {\footnotesize $v_0$};
  \draw[very thick,dashed,-{Latex[length=2.5mm]},blue!60!black] (A) -- ++(1.05,0.15) node[right=1pt] {\footnotesize $v_0'$};
  \draw[-{Latex[length=1.5mm]}] (A)++(0.55,0.1) arc[start angle=6,end angle=31,radius=0.6];
  \node at ($(A)+(0.75,-0.15)$) {\scriptsize $\delta$};
\end{tikzpicture}
\caption{Перенос вектора убеждения $v_0$ вдоль замкнутого геодезического контура $A\to B\to C\to A$ на искривлённом многообразии личности. При возвращении в исходную точку вектор $v_0'$ повёрнут относительно $v_0$ на угол голономии $\delta$, пропорциональный кривизне, охваченной контуром спора.}
\label{fig:holonomy}
\end{figure}

\subsection{Тождество Бьянки для совести}

Введём тензор кривизны характера $R^{\rho}{}_{\sigma\mu\nu}$ и рассмотрим его свёртку -- тензор вины $\mathrm{Ric}_{\mu\nu}$, распределённый по замкнутому контуру социальных отношений (семья, коллектив). Дифференциальное тождество Бьянки, $\nabla_{[\lambda}R_{\rho\sigma]\mu\nu}=0$, применённое к этой системе, даёт следующий результат.

\begin{theorem}[О сохранении вины]
Сколько бы вина ни перераспределялась внутри замкнутого контура отношений, дивергенция суммарного чувства вины по контуру сохраняется тождественно.
\end{theorem}

Иными словами, построить конфликт с нулевым числом виноватых в замкнутой системе структурно невозможно -- вина может быть переопределена, но не аннулирована.

\section{Проблема предвзятости и инвариантная оценка}

\subsection{Координатная зависимость покомпонентных суждений}

Любое суждение вида «он ленивый», «она вспыльчивая», выраженное через отдельные компоненты тензора кривизны характера, координатно-зависимо: оно отражает не столько объект оценки, сколько метрику наблюдателя, деформированную его личной историей и текущим эмоциональным состоянием (см. рис.~\ref{fig:covariance}).

\subsection{Инвариантная оценка: скаляр Кречмана личности}

Единственный класс величин, не меняющих значения при переходе между произвольными системами отсчёта, -- полные свёртки тензора самого с собой. По аналогии со скаляром Кречмана общей теории относительности определим \textbf{инвариантную кривизну личности}:
\[
K_{\text{личн}} = R_{\mu\nu\rho\sigma}\,R^{\mu\nu\rho\sigma}.
\]
Эта величина по построению не зависит от того, из какой точки жизни и в каком настроении наблюдатель смотрит на объект оценки.

\subsection{Проблема потери информации}

Свёртка шестнадцати (и более) независимых компонент тензора Римана в единственное число $K_{\text{личн}}$ -- операция необратимая: она уничтожает всю структуру, оставляя единственное значение вида «инвариантная кривизна личности: 7.2». Результат идеально объективен, абсолютно беспристрастен и практически бесполезен -- что структурно воспроизводит эффект сведения многомерного психологического опросника к четырём буквам типа личности.

\subsection{Локальная инерциальность первого впечатления}

\begin{theorem}[О локальной инерциальности первого впечатления]
В любой точке многообразия личности существует локально-инерциальная система координат, в которой первые производные метрики обращаются в ноль.
\end{theorem}

\begin{figure}[htbp]
\centering
\begin{tikzpicture}[scale=1.1]
  % Искривлённая поверхность (профиль) -- парабола, чтобы касательная считалась точно
  \fill[blue!4] plot[smooth,domain=-3.2:3.2,samples=60]
    ({\x},{0.85 - 0.083*\x*\x}) -- (3.2,-2) -- (-3.2,-2) -- cycle;
  \draw[thick] plot[smooth,domain=-3.2:3.2,samples=60]
    ({\x},{0.85 - 0.083*\x*\x});

  % Точка контакта: P=(1,0.767), наклон кривой в P равен -0.166 (производная -0.166x при x=1)
  \coordinate (P) at (1.0,0.767);
  \fill (P) circle (1.6pt);
  \draw[very thick,red,-{Latex[length=2mm]}] (-0.578,1.029) -- (2.578,0.505);
  \node[red,anchor=west] at (2.63,0.44) {\footnotesize касательная плоскость};

  % Подпись точки через выноску, чтобы не пересекаться с касательной
  \node[align=center,text width=2.6cm] (lbl) at (-1.2,-1.0) {\footnotesize точка первого впечатления ($\approx90$ с)};
  \draw[gray,thin] (lbl.north) -- (P);

  % Лупа
  \draw[gray] (P) circle (0.35);
  \draw[gray] ($(P)+(0.25,-0.25)$) -- ($(P)+(1.3,-1.35)$);
  \begin{scope}[xshift=2.3cm,yshift=-3.0cm]
    \draw[thick] (-1.2,0) -- (1.2,0);
    \draw[very thick,red,-{Latex[length=2mm]}] (-1.2,0.05) -- (1.2,-0.05);
    \fill (0,0) circle (1.6pt);
    \node[below,align=center,text width=3.6cm] at (0,-0.3) {\footnotesize вблизи точки контакта кривизна не ощущается};
  \end{scope}
\end{tikzpicture}
\caption{Локально-инерциальная система координат на многообразии личности. В окрестности точки первого контакта кривизна многообразия ещё не проявляется в наблюдаемых эффектах, и первое впечатление субъективно ощущается как объективное.}
\label{fig:inertial}
\end{figure}

На практике такой системой служат первые девяносто секунд знакомства (рис.~\ref{fig:inertial}): кривизна ещё не успевает проявиться, и первое впечатление ощущается как объективное. Это не свойство человека, а артефакт того, что наблюдатель ещё не отъехал от точки контакта достаточно далеко, чтобы почувствовать кривизну -- как собственных предрассудков, так и предрассудков собеседника.

\section{Тензор непонимания}

Пусть $i$ и $j$ -- два наблюдателя, оценивающих одного и того же человека по общему набору признаков. Введём тензор второго ранга $T_{ij}$, описывающий степень их взаимного расхождения, и разложим его на симметричную и антисимметричную части:
\[
T_{ij} = \underbrace{\tfrac12(T_{ij}+T_{ji})}_{S_{ij}} \;+\; \underbrace{\tfrac12(T_{ij}-T_{ji})}_{A_{ij}}.
\]

\begin{figure}[htbp]
\centering
\begin{tikzpicture}[scale=1.05]
  % --- Симметричная часть ---
  \begin{scope}
    \node[circle,draw,thick,minimum size=0.9cm] (i1) at (0,0) {\footnotesize $i$};
    \node[circle,draw,thick,minimum size=0.9cm] (j1) at (3.2,0) {\footnotesize $j$};
    \draw[very thick,-{Latex[length=2.5mm]}] (i1) to[bend left=18] (j1);
    \draw[very thick,-{Latex[length=2.5mm]}] (j1) to[bend left=18] (i1);
    \node[below] at (1.6,-1.05) {\footnotesize $S_{ij}$: молчаливое взаимное непонимание};
  \end{scope}
  % --- Антисимметричная часть ---
  \begin{scope}[xshift=7.6cm]
    \node[circle,draw,thick,minimum size=0.9cm] (i2) at (0,0) {\footnotesize $i$};
    \node[circle,draw,thick,minimum size=0.9cm] (j2) at (3.2,0) {\footnotesize $j$};
    \draw[very thick,red,-{Latex[length=3mm]}] (i2) to[bend left=35] (j2);
    \draw[very thick,red,{Latex[length=3mm]}-] (i2) to[bend right=35] (j2);
    \node[below] at (1.6,-1.05) {\footnotesize $A_{ij}=-A_{ji}$: вращающий момент спора};
  \end{scope}
\end{tikzpicture}
\caption{Разложение тензора непонимания. Симметричная часть $S_{ij}$ описывает устойчивое взаимное непонимание, стабилизирующее бытовую толерантность. Антисимметричная часть $A_{ij}$ эквивалентна псевдовектору момента и не сокращается при смене ролей наблюдателей -- отсюда «ты меня не понимаешь» / «нет, это ты меня не понимаешь».}
\label{fig:tensor-decomp}
\end{figure}

\textbf{Симметричная часть $S_{ij}$} описывает взаимное и устойчивое непонимание: оба наблюдателя одинаково не улавливают один и тот же аспект объекта и молчаливо соглашаются эту тему не поднимать. Это и есть вещество, из которого построена бытовая толерантность.

\textbf{Антисимметричная часть $A_{ij}$} физически интереснее: как всякий антисимметричный тензор ранга 2, она эквивалентна псевдовектору и несёт момент вращения (рис.~\ref{fig:tensor-decomp}). Поскольку $A_{ij}=-A_{ji}$ по построению, эта часть не сокращается при обмене ролями наблюдателей, а меняет знак.

\section{Усреднение по выборке наблюдателей и миф об объективном мнении}

Естественная надежда на «непредвзятость через демократию мнений» опирается на центральную предельную теорему: усреднение оценок $N$ независимых наблюдателей должно давать дисперсию среднего $\sigma^2/N$. Проблема -- в предположении независимости.

На практике наблюдатели коррелированы: общие друзья читают одни и те же сплетни, семья унаследовала общую оптику, а рекомендательные алгоритмы синхронизируют предвзятость раньше, чем человек успевает сформировать «собственное» мнение. Для наблюдателей с попарной корреляцией $\rho$ дисперсия среднего равна
\[
\mathrm{Var}(\bar X) = \rho\sigma^2 + \frac{(1-\rho)\sigma^2}{N}.
\]
При $N\to\infty$ второе слагаемое исчезает, но первое остаётся.

\begin{theorem}[О нередуцируемом осадке сплетни]
Эффективный размер выборки коррелированных наблюдателей ограничен сверху величиной
\[
n_{\text{эфф}} = \frac{N}{1+(N-1)\rho}
\]
независимо от роста $N$.
\end{theorem}

\begin{figure}[htbp]
\centering
\begin{tikzpicture}
\begin{axis}[
  width=11.5cm, height=7cm,
  xlabel={Число опрошенных наблюдателей $N$},
  ylabel={Эффективный размер выборки $n_{\text{эфф}}$},
  xmin=1, xmax=500,
  ymin=0, ymax=60,
  legend pos=south east,
  legend style={font=\footnotesize},
  grid=major,
  grid style={gray!20},
  label style={font=\footnotesize},
  tick label style={font=\scriptsize},
]
\addplot[very thick,blue,domain=1:500,samples=200] {x/(1+(x-1)*0.02)};
\addlegendentry{$\rho=0.02$ (случайные незнакомцы)}
\addplot[very thick,red,domain=1:500,samples=200] {x/(1+(x-1)*0.10)};
\addlegendentry{$\rho=0.10$ (широкий круг общения)}
\addplot[very thick,orange,domain=1:500,samples=200] {x/(1+(x-1)*0.30)};
\addlegendentry{$\rho=0.30$ (близкий круг друзей)}
\addplot[thick,dashed,gray] coordinates {(1,50)(500,50)};
\end{axis}
\end{tikzpicture}
\caption{Эффективный размер выборки $n_{\text{эфф}}$ как функция числа опрошенных наблюдателей $N$ при разных уровнях корреляции $\rho$. При $\rho=0.30$ (типичный дружеский круг) $n_{\text{эфф}}$ выходит на плато уже при $N\approx30$--$50$ и далее не растёт, сколько бы наблюдателей ни было опрошено.}
\label{fig:neff}
\end{figure}

На практике $n_{\text{эфф}}$ упирается в потолок порядка размера дружеского круга задолго до того, как номинальное число опрошенных превысит сотню (рис.~\ref{fig:neff}). Единственный теоретически корректный способ получить непредвзятый скаляр личности -- усреднять мнения наблюдателей, гарантированно не связанных общими социальными приливными силами ($\rho\to0$), что на практике означает опрашивать людей из другого города, никогда не пересекавшихся с объектом оценки, -- метод, дающий результат, неотличимый от полного отсутствия информации.

\section{Обсуждение и ограничения}

Настоящая модель обладает по меньшей мере тремя структурными ограничениями. Во-первых, инвариантный скаляр $K_{\text{личн}}$ по построению теряет всю информацию о направленности кривизны -- два человека с противоположной, но равной по модулю «проблемностью» характера неразличимы для метода (п.~3.3). Во-вторых, локально-инерциальные системы координат существуют лишь в бесконечно малой окрестности точки контакта; уже вторая встреча требует учёта символов Кристоффеля первого впечатления. В-третьих, метод не устраняет корреляцию наблюдателей, а лишь количественно фиксирует её нередуцируемый вклад (п.~5), что можно рассматривать как единственный строгий результат работы.

\section{Заключение}

Мы показали, что перенос аппарата общей теории относительности на психометрику личности -- метрика, ковариантность, тождества Бьянки, скалярные инварианты, тензорное разложение на симметричную и антисимметричную части, усреднение по коррелированным наблюдателям -- даёт внутренне непротиворечивую, полностью формализуемую и, по-видимому, совершенно неприменимую систему. Мы считаем это достаточным основанием для дальнейшего финансирования.

\section*{Вклад авторов}
Р.\,Г.~Александров: концептуализация, методология, тензорный формализм, написание -- черновик. К.\,С.~Антропикович: программное обеспечение, визуализация, формальный анализ, написание -- рецензирование и редактирование, техническая поддержка бесконечного терпения.

\section*{Финансирование}
Исследование выполнено при поддержке гранта Института Прикладной Метафизики № 2026-$\Omega$-$\infty$ «Кривизна всего».

\section*{Доступность данных}
Данные (компоненты тензора кривизны характера всех участников) доступны по обоснованному запросу -- а также по необоснованному, с той же вероятностью успеха.

\begin{thebibliography}{9}
\bibitem{formalin} Формалин, Ф.~Ф. \textit{О ненулевой кривизне субъективного времени в позднем подростковом возрасте.} Труды Института Прикладной Метафизики, б.г.
\bibitem{kovariant} Ковариантный, Г.~К. \textit{Тождества Бьянки и распределённая вина в замкнутых системах родства.} Неопубликованная рукопись.
\end{thebibliography}

\noindent\textbf{Конфликт интересов.} Не заявлен, хотя, согласно Теореме~1, его отсутствие в замкнутом контуре строго доказать невозможно.

\end{document}
